\subsection{Glossary}
\label{sec:p14-glossary}

This compact logical vocabulary is complemented by the cross-domain lookup
in \cref{sec:p14-symbols-terminology-and-glossary}, which distinguishes
source conventions and overloaded terms across reasoning disciplines.

\begin{description}[style=nextline,leftmargin=0pt]
\item[admissible rule] A rule whose conclusion is derivable after every
substitution making all premises derivable, whether or not the rule itself is
derivable.
\item[algebraic semantics] An interpretation of logical operations and
consequence in a declared class of algebras; see Chapter~\ref{sec:p10-algebraic-and-categorical-logic}.
\item[compactness] The property that semantic consequence from a set is already
consequence from a finite subset; first-order compactness is
Theorem~\ref{thm:fol-compactness}.
\item[complete theory] A consistent theory deciding every sentence of its
language; this differs from completeness of a calculus.
\item[congruence] An operation-compatible equivalence relation on an algebra;
see Definition~\ref{def:p10-algebra}.
\item[conservative extension] An extension proving no new sentences in the old
language beyond those already provable in the base theory.
\item[decidable set] A set with a total computable characteristic function.
\item[derivability] Existence of a derivation in a named calculus; never an
unsubscripted synonym for semantic consequence.
\item[effective procedure] A finitely specified method analyzed through a
declared formal model of computation.
\item[elementary embedding] An embedding preserving every first-order formula
with parameters.
\item[equisatisfiable] Having a model simultaneously in the sense that one
formula is satisfiable iff the other is; weaker than equivalence.
\item[filter] An upward-closed, finite-meet-closed subset of a lattice; see
Definition~\ref{def:p10-boolean-algebra}.
\item[finite validity] Truth in every finite structure of the fixed signature;
see Definition~\ref{def:p09-finite-validity}.
\item[frame] The relational component of a Kripke model, before a valuation is
added.
\item[free algebra] An algebra satisfying the extension universal property from
a chosen generating set.
\item[homomorphism] A map preserving every operation or relation specified by
the context.
\item[interpretation of theories] A translation of one theory into another
with a declared domain formula and preservation of axioms/proofs.
\item[invariant] A property unchanged under a specified class of transformations;
the transformation class is part of the assertion.
\item[Kripke model] A frame together with a valuation or domain/interpretation
data appropriate to the modal language.
\item[Lindenbaum--Tarski algebra] The formula algebra modulo provable
interderivability; see Definition~\ref{def:p10-lindenbaum-algebra}.
\item[model checking] Determining whether a fixed finite structure/state
satisfies a formula; distinct from satisfiability over all structures.
\item[normal form] A formula satisfying a declared syntactic restriction and
connected to the input by equivalence or equisatisfiability as explicitly stated.
\item[proof sketch] An incomplete proof presentation whose omitted steps remain
owed and which is never cited as a complete proof.
\item[realizability] An interpretation assigning computational witnesses to
formulas under a specified applicative structure.
\item[relative consistency] A metatheoretic implication
\(\operatorname{Con}(T)\to\operatorname{Con}(U)\), not an absolute consistency proof.
\item[saturation] Realization of specified types over parameter sets below a
declared cardinal bound.
\item[semantic consequence] Truth preservation over a declared class of
interpretations; see Definition~\ref{def:p11-model-consequence}.
\item[soundness] The assertion that derivability in a named calculus implies
semantic consequence under a named semantics.
\item[stable model] A fixed point of the Gelfond--Lifschitz reduct construction;
see Definition~\ref{def:p09-stable-model}.
\item[Stone space] The compact zero-dimensional Hausdorff dual of a Boolean
algebra; see Definition~\ref{def:p10-stone-space}.
\item[substitution] Capture-avoiding replacement with the exact term/formula
and free-variable conditions stated.
\item[theory] A set of sentences, often deductively closed only when explicitly
said so.
\item[type] A consistent set of formulas in specified free variables over a
declared parameter set.
\item[ultrafilter] A maximal proper filter; see Definition~\ref{def:p10-boolean-algebra}.
\item[validity] Truth in every interpretation in a declared class.
\item[variety] A class of similar algebras axiomatized by equations;
equivalently an HSP-closed class when arbitrary products include the empty
product, by Theorem~\ref{thm:p10-hsp}.
\end{description}
